% unidades/30-repaso-general.tex
\chapter{🎌 総復習 — Repaso General}
\unidadheader{30}{総復習}{Repaso General}{Consolidación de N5 y paso a N4}

\section*{📑 Puntos Clave del Nivel}
En estas 30 unidades hemos cubierto:
\begin{itemize}
  \item \textbf{Fundamentos:} Presentaciones, números, tiempo y existencia.
  \item \textbf{Verbos:} Conjugación en ます, forma て, forma た, forma casual.
  \item \textbf{Adjetivos:} Diferencia entre い-adj y な-adj, sus formas pasadas y negativas.
  \item \textbf{Partículas:} は, が, を, に, へ, で, と, も, より, から, まで.
  \item \textbf{Estructuras N4:} Potencial, pasiva, condicionales (たら, と, ば, なら) y keigo básico.
\end{itemize}

\section*{📈 Siguientes Pasos}
Para continuar hacia el nivel N3, se recomienda:
\begin{enumerate}
  \item \textbf{Kanjis:} Dominar los 300 kanjis básicos (N5 + N4).
  \item \textbf{Lectura:} Empezar a leer textos cortos sin furigana.
  \item \textbf{Escucha:} Escuchar podcasts como "Let's Talk in Japanese".
  \item \textbf{Gramática:} Estudiar formas causativas, nominalización y conectores lógicos complejos.
\end{enumerate}

\begin{tcolorbox}[ejerciciobox, title={🏁 Reto Final}]
Escribe un párrafo de al menos 10 líneas presentándote, hablando de tus gustos, tu rutina diaria, alguna experiencia pasada y tus planes futuros en japonés.
\end{tcolorbox}

\begin{center}
  \vspace{2cm}
  {\Huge \textbf{お疲れ様でした！}} \\[1cm]
  {\Large ¡Felicidades por completar este curso de apuntes!}
\end{center}
